\documentclass[11pt]{article}

% Shared notes settings for Elements of AIML lectures (UPES).
% Keep this file free of \documentclass to allow reuse via \input.

\usepackage[utf8]{inputenc}
\usepackage[T1]{fontenc}
\usepackage[a4paper,margin=1in]{geometry}
\usepackage{hyperref}
\usepackage{xurl}
\usepackage{booktabs}
\usepackage{amsmath}
\usepackage{amssymb}
\usepackage{listings}
\usepackage{xcolor}
\usepackage{enumitem}
\usepackage{parskip}

% Code listings (general-purpose).
\lstset{
  basicstyle=\ttfamily\small,
  keywordstyle=\color{blue},
  commentstyle=\color{gray},
  stringstyle=\color{teal},
  showstringspaces=false,
  columns=fullflexible,
  frame=single,
  framerule=0.2pt,
  breaklines=true
}

\setlist[itemize]{topsep=0.3em,itemsep=0.2em}
\setlist[enumerate]{topsep=0.3em,itemsep=0.2em}

% Simple per-section source line.
\newcommand{\notesources}[1]{\par\noindent{\small\color{gray}\textit{Sources:} \sloppy #1\par}}

% Unit-wise source strings for this subject (used by slides/notes).
% Path is relative to the lecture `latex/` folder (the compilation CWD).
\IfFileExists{../../../_shared/aiml-sources.tex}{%
  % Source strings for Elements of AIML (used by slides + notes).

\newcommand{\AIMLRepoURL}{https://github.com/tali7c/Elements-of-AIML}

\newcommand{\AIMLLabSources}{%
Provost and Fawcett, \textit{Data Science for Business}.;%
McKinney, \textit{Python for Data Analysis}.;%
WEKA Documentation (Waikato Environment for Knowledge Analysis), accessed 2026-02-28.;%
Matplotlib and Seaborn documentation, accessed 2026-02-28.%
}

\newcommand{\AIMLUnitISources}{%
Rich and Knight, \textit{Artificial Intelligence}, McGraw Hill, 2017.;%
Russell and Norvig, \textit{Artificial Intelligence: A Modern Approach} (4e), Ch. 1--2 (Introduction, Intelligent Agents).;%
Poole and Mackworth, \textit{Artificial Intelligence: Foundations of Computational Agents} (2e), selected sections.%
}

\newcommand{\AIMLUnitIISources}{%
Russell and Norvig, \textit{Artificial Intelligence: A Modern Approach} (4e), Ch. 7--9 (Logical Agents, First-Order Logic, Inference).;%
Huth and Ryan, \textit{Logic in Computer Science} (2e), selected sections on propositional and first-order logic.;%
Genesereth and Nilsson, \textit{Logical Foundations of Artificial Intelligence}, selected chapters.%
}

\newcommand{\AIMLUnitIIISources}{%
Mueller and Massaron, \textit{Machine Learning for Dummies}, For Dummies, 2016.;%
Mitchell, \textit{Machine Learning}, selected chapters on learning basics and evaluation.;%
Scikit-learn documentation, accessed 2026-02-28.%
}

\newcommand{\AIMLUnitIVSources}{%
Mueller and Massaron, \textit{Machine Learning for Dummies}, For Dummies, 2016.;%
James, Witten, Hastie, and Tibshirani, \textit{An Introduction to Statistical Learning} (2e), selected chapters.;%
Scikit-learn documentation, accessed 2026-02-28.%
}

\newcommand{\AIMLUnitVSources}{%
NIST, \textit{AI Risk Management Framework (AI RMF 1.0)}, 2023.;%
Barocas, Hardt, and Narayanan, \textit{Fairness and Machine Learning} (online book), selected sections.;%
Knaflic, \textit{Storytelling with Data}, selected chapters on communicating results.%
}
%
}{%
  % Faculty copies live under `private/` so their compile CWD is different.
  \IfFileExists{../../../../public/_shared/aiml-sources.tex}{%
    % Source strings for Elements of AIML (used by slides + notes).

\newcommand{\AIMLRepoURL}{https://github.com/tali7c/Elements-of-AIML}

\newcommand{\AIMLLabSources}{%
Provost and Fawcett, \textit{Data Science for Business}.;%
McKinney, \textit{Python for Data Analysis}.;%
WEKA Documentation (Waikato Environment for Knowledge Analysis), accessed 2026-02-28.;%
Matplotlib and Seaborn documentation, accessed 2026-02-28.%
}

\newcommand{\AIMLUnitISources}{%
Rich and Knight, \textit{Artificial Intelligence}, McGraw Hill, 2017.;%
Russell and Norvig, \textit{Artificial Intelligence: A Modern Approach} (4e), Ch. 1--2 (Introduction, Intelligent Agents).;%
Poole and Mackworth, \textit{Artificial Intelligence: Foundations of Computational Agents} (2e), selected sections.%
}

\newcommand{\AIMLUnitIISources}{%
Russell and Norvig, \textit{Artificial Intelligence: A Modern Approach} (4e), Ch. 7--9 (Logical Agents, First-Order Logic, Inference).;%
Huth and Ryan, \textit{Logic in Computer Science} (2e), selected sections on propositional and first-order logic.;%
Genesereth and Nilsson, \textit{Logical Foundations of Artificial Intelligence}, selected chapters.%
}

\newcommand{\AIMLUnitIIISources}{%
Mueller and Massaron, \textit{Machine Learning for Dummies}, For Dummies, 2016.;%
Mitchell, \textit{Machine Learning}, selected chapters on learning basics and evaluation.;%
Scikit-learn documentation, accessed 2026-02-28.%
}

\newcommand{\AIMLUnitIVSources}{%
Mueller and Massaron, \textit{Machine Learning for Dummies}, For Dummies, 2016.;%
James, Witten, Hastie, and Tibshirani, \textit{An Introduction to Statistical Learning} (2e), selected chapters.;%
Scikit-learn documentation, accessed 2026-02-28.%
}

\newcommand{\AIMLUnitVSources}{%
NIST, \textit{AI Risk Management Framework (AI RMF 1.0)}, 2023.;%
Barocas, Hardt, and Narayanan, \textit{Fairness and Machine Learning} (online book), selected sections.;%
Knaflic, \textit{Storytelling with Data}, selected chapters on communicating results.%
}
%
  }{%
    % Source strings for Elements of AIML (used by slides + notes).

\newcommand{\AIMLRepoURL}{https://github.com/tali7c/Elements-of-AIML}

\newcommand{\AIMLLabSources}{%
Provost and Fawcett, \textit{Data Science for Business}.;%
McKinney, \textit{Python for Data Analysis}.;%
WEKA Documentation (Waikato Environment for Knowledge Analysis), accessed 2026-02-28.;%
Matplotlib and Seaborn documentation, accessed 2026-02-28.%
}

\newcommand{\AIMLUnitISources}{%
Rich and Knight, \textit{Artificial Intelligence}, McGraw Hill, 2017.;%
Russell and Norvig, \textit{Artificial Intelligence: A Modern Approach} (4e), Ch. 1--2 (Introduction, Intelligent Agents).;%
Poole and Mackworth, \textit{Artificial Intelligence: Foundations of Computational Agents} (2e), selected sections.%
}

\newcommand{\AIMLUnitIISources}{%
Russell and Norvig, \textit{Artificial Intelligence: A Modern Approach} (4e), Ch. 7--9 (Logical Agents, First-Order Logic, Inference).;%
Huth and Ryan, \textit{Logic in Computer Science} (2e), selected sections on propositional and first-order logic.;%
Genesereth and Nilsson, \textit{Logical Foundations of Artificial Intelligence}, selected chapters.%
}

\newcommand{\AIMLUnitIIISources}{%
Mueller and Massaron, \textit{Machine Learning for Dummies}, For Dummies, 2016.;%
Mitchell, \textit{Machine Learning}, selected chapters on learning basics and evaluation.;%
Scikit-learn documentation, accessed 2026-02-28.%
}

\newcommand{\AIMLUnitIVSources}{%
Mueller and Massaron, \textit{Machine Learning for Dummies}, For Dummies, 2016.;%
James, Witten, Hastie, and Tibshirani, \textit{An Introduction to Statistical Learning} (2e), selected chapters.;%
Scikit-learn documentation, accessed 2026-02-28.%
}

\newcommand{\AIMLUnitVSources}{%
NIST, \textit{AI Risk Management Framework (AI RMF 1.0)}, 2023.;%
Barocas, Hardt, and Narayanan, \textit{Fairness and Machine Learning} (online book), selected sections.;%
Knaflic, \textit{Storytelling with Data}, selected chapters on communicating results.%
}
%
  }%
}


\title{Elements of AIML Lab\\Experiment 04 (After-submission Reference): Pandas Import/Export and Summaries}
\author{Tofik Ali}
\date{\today}

\begin{document}
\maketitle
\tableofcontents

\section{How to use this reference}
This reference is meant to be shared \textbf{after you submit} Experiment~04.
It provides a complete, readable walkthrough of the \textbf{import $\rightarrow$ inspect $\rightarrow$ clean $\rightarrow$ export} cycle.

\section{What this experiment is really testing}
Before we build ML models, we must be able to answer:
\begin{itemize}[leftmargin=*]
  \item What are the columns and data types?
  \item How many rows/columns do we have?
  \item Are there missing values, duplicates, impossible values?
  \item Can we export a processed dataset for later experiments?
\end{itemize}

\section{Worked example (self-contained)}
\subsection{Step 0: create a small example CSV (only for demonstration)}
If you already used a real dataset, skip to Step~1. The code below creates an example file.
\begin{lstlisting}[language=Python]
import pandas as pd

df_demo = pd.DataFrame({
    "age": [21, 22, None, 24, 24],
    "city": ["Dehradun", "Delhi", "Delhi", None, "Mumbai"],
    "marks": [78, 85, 66, 92, 92],
})
df_demo.to_csv("demo.csv", index=False)
\end{lstlisting}

\subsection{Step 1: import the dataset}
\begin{lstlisting}[language=Python]
import pandas as pd

df = pd.read_csv("demo.csv")   # replace with your dataset path
\end{lstlisting}
If your CSV uses a different separator, try \texttt{sep=";"} (or check the file).

\subsection{Step 2: basic inspection (what each command tells you)}
\begin{lstlisting}[language=Python]
print("Shape:", df.shape)     # (rows, columns)
print(df.head())              # preview
print(df.info())              # dtypes + non-null counts
\end{lstlisting}
\textbf{How to interpret:}
\begin{itemize}[leftmargin=*]
  \item \texttt{shape}: helps you confirm import success (not empty, not truncated).
  \item \texttt{head}: checks if columns look correct (no shifted columns due to commas).
  \item \texttt{info}: shows types; if a numeric column appears as \texttt{object}, it may contain strings like ``--''.
\end{itemize}

\subsection{Step 3: missing values and duplicates}
\begin{lstlisting}[language=Python]
missing = df.isna().sum().sort_values(ascending=False)
print("Missing per column:\n", missing)
print("Duplicates:", df.duplicated().sum())
\end{lstlisting}
\textbf{Interpretation:} columns with high missingness may require imputation or removal in later experiments.

\subsection{Step 4: descriptive summaries}
\begin{lstlisting}[language=Python]
print(df.describe())                  # numeric columns
print(df.describe(include="object"))  # categorical columns
\end{lstlisting}
Look for:
\begin{itemize}[leftmargin=*]
  \item suspicious min/max (data entry errors/outliers),
  \item impossible values (negative age),
  \item rare categories (very low counts).
\end{itemize}

\subsection{Step 5: a small clean step + export}
One realistic mini-clean is: drop duplicates, fill a missing category with \texttt{"Unknown"}, and export.
\begin{lstlisting}[language=Python]
df2 = df.drop_duplicates()
df2["city"] = df2["city"].fillna("Unknown")
df2.to_csv("processed.csv", index=False)
\end{lstlisting}
\textbf{Cross-check:} Open \texttt{processed.csv} and confirm it contains the expected rows/columns.

\section{What to write in the report (example phrasing)}
\begin{itemize}[leftmargin=*]
  \item ``The dataset has 1200 rows and 14 columns. The column \texttt{Age} has 23 missing values (1.9\%).''
  \item ``\texttt{info()} shows \texttt{Income} is \texttt{object}; we found currency symbols and removed them (converted to float).''
  \item ``We exported \texttt{processed.csv} after dropping duplicates and filling missing city values with \texttt{Unknown}.''
\end{itemize}

\section{Troubleshooting (common issues and fixes)}
\begin{itemize}[leftmargin=*]
  \item \textbf{Unicode errors:} try \texttt{pd.read\_csv(..., encoding="utf-8")} or \texttt{"latin-1"}.
  \item \textbf{Wrong separator:} use \texttt{sep=";"} or \texttt{sep="\\t"} for TSV.
  \item \textbf{Numeric column read as text:} use \texttt{pd.to\_numeric(..., errors="coerce")}.
  \item \textbf{Dates read as strings:} use \texttt{parse\_dates=[...]} or \texttt{pd.to\_datetime}.
\end{itemize}

\notesources{\AIMLLabSources}
\end{document}

